\documentclass[a4paper,12pt,twoside]{memoir}

% Castellano
\usepackage[spanish,es-tabla]{babel}
\selectlanguage{spanish}
\usepackage[utf8]{inputenc}
\usepackage[T1]{fontenc}
\usepackage{lmodern} % Scalable font
\usepackage{microtype}
\usepackage{placeins}

\RequirePackage{booktabs}
\RequirePackage[table]{xcolor}
\RequirePackage{xtab}
\RequirePackage{multirow}

% Links
\usepackage[colorlinks]{hyperref}
\hypersetup{
	allcolors = {red}
}

% Ecuaciones
\usepackage{amsmath}

% Rutas de fichero / paquete
\newcommand{\ruta}[1]{{\sffamily #1}}

% Párrafos
\nonzeroparskip


% Imagenes
\usepackage{graphicx}
\newcommand{\imagen}[2]{
	\begin{figure}[!h]
		\centering
		\includegraphics[width=0.9\textwidth]{#1}
		\caption{#2}\label{fig:#1}
	\end{figure}
	\FloatBarrier
}

\newcommand{\imagenflotante}[2]{
	\begin{figure}%[!h]
		\centering
		\includegraphics[width=0.9\textwidth]{#1}
		\caption{#2}\label{fig:#1}
	\end{figure}
}



% El comando \figura nos permite insertar figuras comodamente, y utilizando
% siempre el mismo formato. Los parametros son:
% 1 -> Porcentaje del ancho de página que ocupará la figura (de 0 a 1)
% 2 --> Fichero de la imagen
% 3 --> Texto a pie de imagen
% 4 --> Etiqueta (label) para referencias
% 5 --> Opciones que queramos pasarle al \includegraphics
% 6 --> Opciones de posicionamiento a pasarle a \begin{figure}
\newcommand{\figuraConPosicion}[6]{%
  \setlength{\anchoFloat}{#1\textwidth}%
  \addtolength{\anchoFloat}{-4\fboxsep}%
  \setlength{\anchoFigura}{\anchoFloat}%
  \begin{figure}[#6]
    \begin{center}%
      \Ovalbox{%
        \begin{minipage}{\anchoFloat}%
          \begin{center}%
            \includegraphics[width=\anchoFigura,#5]{#2}%
            \caption{#3}%
            \label{#4}%
          \end{center}%
        \end{minipage}
      }%
    \end{center}%
  \end{figure}%
}

%
% Comando para incluir imágenes en formato apaisado (sin marco).
\newcommand{\figuraApaisadaSinMarco}[5]{%
  \begin{figure}%
    \begin{center}%
    \includegraphics[angle=90,height=#1\textheight,#5]{#2}%
    \caption{#3}%
    \label{#4}%
    \end{center}%
  \end{figure}%
}
% Para las tablas
\newcommand{\otoprule}{\midrule [\heavyrulewidth]}
%
% Nuevo comando para tablas pequeñas (menos de una página).
\newcommand{\tablaSmall}[5]{%
 \begin{table}
  \begin{center}
   \rowcolors {2}{gray!35}{}
   \begin{tabular}{#2}
    \toprule
    #4
    \otoprule
    #5
    \bottomrule
   \end{tabular}
   \caption{#1}
   \label{tabla:#3}
  \end{center}
 \end{table}
}

%
% Nuevo comando para tablas pequeñas (menos de una página).
\newcommand{\tablaSmallSinColores}[5]{%
 \begin{table}[H]
  \begin{center}
   \begin{tabular}{#2}
    \toprule
    #4
    \otoprule
    #5
    \bottomrule
   \end{tabular}
   \caption{#1}
   \label{tabla:#3}
  \end{center}
 \end{table}
}

\newcommand{\tablaApaisadaSmall}[5]{%
\begin{landscape}
  \begin{table}
   \begin{center}
    \rowcolors {2}{gray!35}{}
    \begin{tabular}{#2}
     \toprule
     #4
     \otoprule
     #5
     \bottomrule
    \end{tabular}
    \caption{#1}
    \label{tabla:#3}
   \end{center}
  \end{table}
\end{landscape}
}

%
% Nuevo comando para tablas grandes con cabecera y filas alternas coloreadas en gris.
\newcommand{\tabla}[6]{%
  \begin{center}
    \tablefirsthead{
      \toprule
      #5
      \otoprule
    }
    \tablehead{
      \multicolumn{#3}{l}{\small\sl continúa desde la página anterior}\\
      \toprule
      #5
      \otoprule
    }
    \tabletail{
      \hline
      \multicolumn{#3}{r}{\small\sl continúa en la página siguiente}\\
    }
    \tablelasttail{
      \hline
    }
    \bottomcaption{#1}
    \rowcolors {2}{gray!35}{}
    \begin{xtabular}{#2}
      #6
      \bottomrule
    \end{xtabular}
    \label{tabla:#4}
  \end{center}
}

%
% Nuevo comando para tablas grandes con cabecera.
\newcommand{\tablaSinColores}[6]{%
  \begin{center}
    \tablefirsthead{
      \toprule
      #5
      \otoprule
    }
    \tablehead{
      \multicolumn{#3}{l}{\small\sl continúa desde la página anterior}\\
      \toprule
      #5
      \otoprule
    }
    \tabletail{
      \hline
      \multicolumn{#3}{r}{\small\sl continúa en la página siguiente}\\
    }
    \tablelasttail{
      \hline
    }
    \bottomcaption{#1}
    \begin{xtabular}{#2}
      #6
      \bottomrule
    \end{xtabular}
    \label{tabla:#4}
  \end{center}
}

%
% Nuevo comando para tablas grandes sin cabecera.
\newcommand{\tablaSinCabecera}[5]{%
  \begin{center}
    \tablefirsthead{
      \toprule
    }
    \tablehead{
      \multicolumn{#3}{l}{\small\sl continúa desde la página anterior}\\
      \hline
    }
    \tabletail{
      \hline
      \multicolumn{#3}{r}{\small\sl continúa en la página siguiente}\\
    }
    \tablelasttail{
      \hline
    }
    \bottomcaption{#1}
  \begin{xtabular}{#2}
    #5
   \bottomrule
  \end{xtabular}
  \label{tabla:#4}
  \end{center}
}



\definecolor{cgoLight}{HTML}{EEEEEE}
\definecolor{cgoExtralight}{HTML}{FFFFFF}

%
% Nuevo comando para tablas grandes sin cabecera.
\newcommand{\tablaSinCabeceraConBandas}[5]{%
  \begin{center}
    \tablefirsthead{
      \toprule
    }
    \tablehead{
      \multicolumn{#3}{l}{\small\sl continúa desde la página anterior}\\
      \hline
    }
    \tabletail{
      \hline
      \multicolumn{#3}{r}{\small\sl continúa en la página siguiente}\\
    }
    \tablelasttail{
      \hline
    }
    \bottomcaption{#1}
    \rowcolors[]{1}{cgoExtralight}{cgoLight}

  \begin{xtabular}{#2}
    #5
   \bottomrule
  \end{xtabular}
  \label{tabla:#4}
  \end{center}
}


















\graphicspath{ {./img/} }

% Capítulos
\chapterstyle{bianchi}
\newcommand{\capitulo}[2]{
	\setcounter{chapter}{#1}
	\setcounter{section}{0}
	\chapter*{#2}
	\addcontentsline{toc}{chapter}{#2}
	\markboth{#2}{#2}
}

% Apéndices
\renewcommand{\appendixname}{Apéndice}
\renewcommand*\cftappendixname{\appendixname}

\newcommand{\apendice}[1]{
	%\renewcommand{\thechapter}{A}
	\chapter{#1}
}

\renewcommand*\cftappendixname{\appendixname\ }

% Formato de portada
\makeatletter
\usepackage{xcolor}
\newcommand{\tutor}[1]{\def\@tutor{#1}}
\newcommand{\course}[1]{\def\@course{#1}}
\definecolor{cpardoBox}{HTML}{E6E6FF}
\def\maketitle{
  \null
  \thispagestyle{empty}
  % Cabecera ----------------
\noindent\includegraphics[width=\textwidth]{cabecera}\vspace{1cm}%
  \vfill
  % Título proyecto y escudo informática ----------------
  \colorbox{cpardoBox}{%
    \begin{minipage}{.8\textwidth}
      \vspace{.5cm}\Large
      \begin{center}
      \textbf{TFG del Grado en Ingeniería Informática}\vspace{.6cm}\\
      \textbf{\LARGE\@title{}}
      \end{center}
      \vspace{.2cm}
    \end{minipage}

  }%
  \hfill\begin{minipage}{.20\textwidth}
    \includegraphics[width=\textwidth]{escudoInfor}
  \end{minipage}
  \vfill
  % Datos de alumno, curso y tutores ------------------
  \begin{center}%
  {%
    \noindent\LARGE
    Presentado por \@author{}\\ 
    en Universidad de Burgos --- \@date{}\\
    Tutor: \@tutor{}\\
  }%
  \end{center}%
  \null
  \cleardoublepage
  }
\makeatother

\newcommand{\nombre}{Nombre del alumno} %%% cambio de comando

% Datos de portada
\title{título del TFG}
\author{\nombre}
\tutor{nombre tutor}
\date{\today}

\begin{document}

\maketitle


\newpage\null\thispagestyle{empty}\newpage


%%%%%%%%%%%%%%%%%%%%%%%%%%%%%%%%%%%%%%%%%%%%%%%%%%%%%%%%%%%%%%%%%%%%%%%%%%%%%%%%%%%%%%%%
\thispagestyle{empty}


\noindent\includegraphics[width=\textwidth]{cabecera}\vspace{1cm}

\noindent D. nombre tutor, profesor del departamento de nombre departamento, área de nombre área.

\noindent Expone:

\noindent Que el alumno D. \nombre, con DNI dni, ha realizado el Trabajo final de Grado en Ingeniería Informática titulado título de TFG. 

\noindent Y que dicho trabajo ha sido realizado por el alumno bajo la dirección del que suscribe, en virtud de lo cual se autoriza su presentación y defensa.

\begin{center} %\large
En Burgos, {\large \today}
\end{center}

\vfill\vfill\vfill

% Author and supervisor
\begin{minipage}{0.45\textwidth}
\begin{flushleft} %\large
Vº. Bº. del Tutor:\\[2cm]
D. nombre tutor
\end{flushleft}
\end{minipage}
\hfill
\begin{minipage}{0.45\textwidth}
\begin{flushleft} %\large
Vº. Bº. del co-tutor:\\[2cm]
D. nombre co-tutor
\end{flushleft}
\end{minipage}
\hfill

\vfill

% para casos con solo un tutor comentar lo anterior
% y descomentar lo siguiente
%Vº. Bº. del Tutor:\\[2cm]
%D. nombre tutor


\newpage\null\thispagestyle{empty}\newpage




\frontmatter

% Abstract en castellano
\renewcommand*\abstractname{Resumen}
\begin{abstract}
En este primer apartado se hace una \textbf{breve} presentación del tema que se aborda en el proyecto.
\end{abstract}

\renewcommand*\abstractname{Descriptores}
\begin{abstract}
Palabras separadas por comas que identifiquen el contenido del proyecto Ej: servidor web, buscador de vuelos, android \ldots
\end{abstract}

\clearpage

% Abstract en inglés
\renewcommand*\abstractname{Abstract}
\begin{abstract}
A \textbf{brief} presentation of the topic addressed in the project.
\end{abstract}

\renewcommand*\abstractname{Keywords}
\begin{abstract}
keywords separated by commas.
\end{abstract}

\clearpage

% Indices
\tableofcontents

\clearpage

\listoffigures

\clearpage

\listoftables
\clearpage

\mainmatter
\capitulo{1}{Introducción}

Descripción del contenido del trabajo y del estrucutra de la memoria y del resto de materiales entregados.

\capitulo{2}{Objetivos del proyecto}

Los objetivos que se muestran a continuación están basados en un punto principal, que es la mejora de la aplicación ya existente, añadiendo nuevas funcionalidades que la hagan más atractiva para el usuario.

Se van a llevar a cabo las siguientes mejoras:
\begin{itemize}
\tightlist
\item Actualizar la base de datos con información referente a horarios de apertura y cierre de los distintos puntos de interés.
\item Actualizar las etiquetas de los diferentes puntos de interés para poder buscar por filtros personalizados.
\item Implementar un algoritmo que tenga en cuenta los horarios de apertura y cierre de los distintos puntos de interés, para añadir realismo a la aplicación y generar rutas más eficientes, ya que el modelo anterior, al no tener en cuenta estos datos, podía añadir a la ruta un punto de interés "no visitable".
\item Implementar la funcionalidad de no repetir puntos de interés en caso de estar planificando la ruta en una ciudad ya visitada.
\end{itemize}

Los objetivos personales son:
\begin{itemize}
\tightlist
\item Aprender a utilizar las herramientas PostgreSQL, PostGIS y GlassFish para tener una base sólida en el campo de creación y uso de servidores.
\item Aprender a utilizar el sistema de composición de textos LaTeX, así como la herramienta de desarrollo Texmaker.
\item Aprender a trabajar en un "equipo de desarrollo" y mejorar en el mismo para tener una experiencia más parecida a la que podemos encontrarnos a la hora de realizar un proyecto "real".
\end{itemize}
\capitulo{3}{Conceptos teóricos}

%En aquellos proyectos que necesiten para su comprensión y desarrollo de unos conceptos teóricos de una determinada materia o de un determinado dominio de conocimiento, debe existir un apartado que sintetice dichos conceptos.
%
%Algunos conceptos teóricos de \LaTeX \footnote{Créditos a los proyectos de Álvaro López Cantero: Configurador de Presupuestos y Roberto Izquierdo Amo: PLQuiz}.
%
%\section{Secciones}
%
%Las secciones se incluyen con el comando section.
%
%\subsection{Subsecciones}
%
%Además de secciones tenemos subsecciones.
%
%\subsubsection{Subsubsecciones}
%
%Y subsecciones. 
%
%
%\section{Referencias}
%
%Las referencias se incluyen en el texto usando cite \cite{wiki:latex}. Para citar webs, artículos o libros \cite{koza92}.
%
%
%\section{Imágenes}
%
%Se pueden incluir imágenes con los comandos standard de \LaTeX, pero esta plantilla dispone de comandos %propios como por ejemplo el siguiente:
%
%\imagen{escudoInfor}{Autómata para una expresión vacía}
%
%
%
%\section{Listas de items}
%
%Existen tres posibilidades:
%
%\begin{itemize}
%	\item primer item.
%	\item segundo item.
%\end{itemize}
%
%\begin{enumerate}
%	\item primer item.
%	\item segundo item.
%\end{enumerate}
%
%\begin{description}
%	\item[Primer item] más información sobre el primer item.
%	\item[Segundo item] más información sobre el segundo item.
%\end{description}
%	
%\begin{itemize}
%\item 
%\end{itemize}
%
%\section{Tablas}
%
%Igualmente se pueden usar los comandos específicos de \LaTeX o bien usar alguno de los comandos de la plantilla.
%
%\tablaSmall{Herramientas y tecnologías utilizadas en cada parte del proyecto}{l c c c c}{herramientasportipodeuso}
%{ \multicolumn{1}{l}{Herramientas} & App AngularJS & API REST & BD & Memoria \\}{ 
%HTML5 & X & & &\\
%CSS3 & X & & &\\
%BOOTSTRAP & X & & &\\
%JavaScript & X & & &\\
%AngularJS & X & & &\\
%Bower & X & & &\\
%PHP & & X & &\\
%Karma + Jasmine & X & & &\\
%Slim framework & & X & &\\
%Idiorm & & X & &\\
%Composer & & X & &\\
%JSON & X & X & &\\
%PhpStorm & X & X & &\\
%MySQL & & & X &\\
%PhpMyAdmin & & & X &\\
%Git + BitBucket & X & X & X & X\\
%Mik\TeX{} & & & & X\\
%\TeX{}Maker & & & & X\\
%Astah & & & & X\\
%Balsamiq Mockups & X & & &\\
%VersionOne & X & X & X & X\\
%} 

En este capítulo se describirán, como ya anticipa el título, diversos aspectos relacionados con cuestiones teóricas, que permitirán al lector una correcta comprensión del trabajo.

\section{The Orienteering Problem}
%Sección OP
El ``problema de la orientación'' u ``Orienteering Problem'' (OP)\cite{op} propone una situación con diferentes puntos que se pueden visitar, cada uno con una puntuación o score asociado.
El objetivo es, partiendo de un punto determinado, llegar al destino dentro de un límite de tiempo maximizando la puntuación total obtenida de visitar los diferentes puntos propuestos. Cada punto disponible podrá visitarse como máximo una sola vez.

Aunque se pueden ver ciertas similitudes con el ``Problema del Viajante'' (``Travelling Salesman Problem'')\cite{wiki:TSP} debemos considerarlo como un problema totalmente distinto, debido a que NO es necesario visitar todos los nodos disponibles, debido al límite de tiempo.
Podemos ver al ``Problema de la Orientación'' como una combinación entre el ``Problema del Viajante'' y el ``Problema de la Mochila''\cite{wiki:KP}.

De este problema original han surgido diferentes variantes a medida que se han ido proponiendo nuevos escenarios, como pueden ser \textit{trabajando en equipo (TOP)}, teniendo en cuenta \textit{ventanas de tiempo (OPTW)}, \textit{dependiente del tiempo (TDOP)} y todas las combinaciones posibles de los anteriores, por lo que podemos encontrarnos con el caso de ``Problema de la orientación por equipos dependiente del tiempo y con ventanas de tiempo" (TDTOPTW).
\subsection{The Orienteering Problem with Time Windows}
%Subsección OPTW
El ``problema de la orientación con ventanas de tiempo'' (OPTW)\cite{optw} considera unos márgenes de tiempo para cada punto disponible para visitar. La visita a un punto solo puede comenzar dentro de sus márgenes de tiempo [O{i},C{i}].

En caso de llegar antes del horario de apertura (O{i}) será necesario esperar. Si por el contrario llegamos después del horario de cierre(C{i}) la visita a ese punto será inviable y por tanto no podrá comenzar (no se obtendrá recompensa).

\section{Sistema de información geográfica}
%Sección datos GIS
Un sistema de información geográfica (GIS) \cite{wiki:gis} es un conjunto de herramientas que integra y relaciona diversos componentes que permiten la organización, almacenamiento, manipulación, análisis y modelización de grandes cantidades de datos procedentes del mundo real que están vinculados a una referencia espacial, facilitando la incorporación de aspectos sociales-culturales, económicos y ambientales que conducen a la toma de decisiones de una manera más eficaz.

Funcionalidad de los GIS:
\begin{enumerate}
\item Almacenar información, que se puede obtener por diferentes métodos, como pueden ser: GPS, fotografía aérea, imágenes satélite, bases de datos \ldots
\item Visualizar datos almacenados.
\item Hacer consulta sobre los datos seleccionados. De esta forma podemos presentar la información de una forma mucho más amigable para el usuario, mediante mapas, gráficos, tablas \ldots
\item Hacer análisis para generar nuevas capas de información útil.
\end{enumerate}

Los tipos de datos geográficos más utilizados actualmente por las GIS son:
\begin{itemize}
\item Datos Vectoriales, utilizados para representar fenómenos discretos, describiendo objetos geográficos a partir de vectores definidos por pares de coordenadas.
\item Datos Raster, utilizdos para representar fenómenos no discretos. Divide el área a representar en una retícula de celdas y atribuye a cada una de estas un valor numérico para representar su valor temático. Cabe destacar que el origen siempre esta en la esquina superior izquierda de la retícula.
\end{itemize}

\section{Scrum}
\textit{Scrum} es un marco de trabajo para desarrollo ágil de software.\cite{wiki:scrum}
Es una metodología muy popular en la actualidad y aunque su enfoque inicial es para proyectos de software, puede ser fácilmente adaptable a otros contextos.

Si no se siguen bien sus reglas puede ser algo complicado de implementar, pero en general es una metodología muy fácil de comprender.

\subsection{Scrum Team}
\textit{Scrum} se compone de tres roles principales:
\begin{itemize}
\item \textbf{Product Owner: }Es la persona que representa al cliente y tiene la misión de conocer todas las necesidades de este. Deberá transmitir esta información tanto al \textit{Scrum Master} como al \textit{Development Team}.
\item \textbf{Scrum Master: }Es el moderador del grupo de trabajo. Aunque esta posición puede guardar similitudes con la de un líder, el \textit{Scrum Master} no es el encargado ni de dar órdenes ni de decidir el modo de hacer las cosas. Su función es asegurarse de que el \textit{Development Team} entienda y trabaje en las necesidades del cliente
\item \textbf{Development Team: }Son las personas capacitadas para crear la solución que el cliente necesita.
\end{itemize}

\subsection{Ciclo de Vida Scrum}
Inicialmente el \textit{Product Owner} definirá un artefacto llamado \textit{Product Backlog}, donde plasmará todas las solicitudes del cliente y será presentado al \textit{Scrum Team} en una reunión llamada \textit{Sprint Planning Meeting}.
Como resultado de esta primera reunión se obtendrá una lista de funcionalidades tomadas del \textit{Product Backlog} a las que se asigna el nombre de \textit{Sprint Backlog}.
\imagen{scrumm}{Ciclo de Vida Scrum\cite{wiki:scrum}}
Las funcionalidades definidas en el \textit{Sprint Backlog} deberán completarse en un periodo de 1 a 4 semanas denominado \textit{Sprint}.

\subsection{Sprint}
Es el corazón del \textit{Scrum}, ya que responde al proceso de construcción de las necesidades del cliente, pero divididas en un módulo funcional. 
El tiempo de desarrollo debe ser entre 1 y 4 semanas, dependiendo de como de complejas sean las funcionalidades acordadas en el \textit{Sprint Backlog}.

Las partes que participan en el \textit{Sprint} son el \textit{Scrum Master} y el \textit{Development Team}, siendo este último el encargado de construir la necesidad que construye el \textit{Sprint}. El primero tendrá la misión de ayudar al equipo de desarrollo.

\subsection{Las reuniones}
Las reuniones son uno de los elementos característicos del \textit{Scrum}. Estas reuniones tienen lugar en cada uno de los \textit{Sprints} y son las siguientes:
\subsubsection{Sprint Planning Meeting}
Como se ha mencionado previamente es el primer paso al comenzar un \textit{Sprint}.Se definirá un \textit{Sprint Backlog} con las funcionalidades a completar en el \textit{Sprint} actual.
\subsubsection{DailyScrum}
Es una de las actividades representativas del \textit{Scrum} y tiene como finalidad hacer un seguimiento diario del proceso de desarrollo.

Consiste en una reunión entre el \textit{Scrum Master} y el \textit{Development Team}, en la que se harán una serie de preguntas muy puntuales a cada pregunta dentro del equipo de desarrollo. Estas preguntas son: qué se hizo ayer, qué se va a hacer hoy, qué se va a hacer mañana y que problemas se han encontrado.

Debido a la naturaleza informativa de esta reunión, su duración deberá ser entre 5 y 15 minutos. 
\subsubsection{Sprint Review}
Se trata de una de las dos reuniones que tienen lugar al finalizar un \textit{Sprint}. En esta estarán involucrados tanto el \textit{Scrum Master} como el \textit{Product Owner} y \textit{Development Team} para verificar el cumplimiento de los objetivos de ese \textit{Sprint} y así garantizar los tiempos de entrega del producto.
\subsubsection{Retrospectiva del Sprint}
La segunda de las reuniones que tiene lugar después de cada \textit{Sprint}. Se analizarán los resultados del \textit{Sprint} recién acabado, para poder solventar problemáticas y mejorar el proceso para \textit{Sprints} posteriores. Nada más acabar un Sprint se iniciará uno nuevo, volviendo a repetir el ciclo.

\capitulo{4}{Técnicas y herramientas}

Esta parte de la memoria tiene como objetivo presentar las técnicas metodológicas y las herramientas de desarrollo que se han utilizado para llevar a cabo el proyecto. Si se han estudiado diferentes alternativas de metodologías, herramientas, bibliotecas se puede hacer un resumen de los aspectos más destacados de cada alternativa, incluyendo comparativas entre las distintas opciones y una justificación de las elecciones realizadas. 
No se pretende que este apartado se convierta en un capítulo de un libro dedicado a cada una de las alternativas, sino comentar los aspectos más destacados de cada opción, con un repaso somero a los fundamentos esenciales y referencias bibliográficas para que el lector pueda ampliar su conocimiento sobre el tema.



\capitulo{5}{Aspectos relevantes del desarrollo del proyecto}

%Este apartado pretende recoger los aspectos más interesantes del desarrollo del proyecto, comentados por los autores del mismo.
%Debe incluir desde la exposición del ciclo de vida utilizado, hasta los detalles de mayor relevancia de las fases de análisis, diseño e implementación.
%Se busca que no sea una mera operación de copiar y pegar diagramas y extractos del código fuente, sino que realmente se justifiquen los caminos de solución que se han tomado, especialmente aquellos que no sean triviales.
%Puede ser el lugar más adecuado para documentar los aspectos más interesantes del diseño y de la implementación, con un mayor hincapié en aspectos tales como el tipo de arquitectura elegido, los índices de las tablas de la base de datos, normalización y desnormalización, distribución en ficheros3, reglas de negocio dentro de las bases de datos (EDVHV GH GDWRV DFWLYDV), aspectos de desarrollo relacionados con el WWW...
%Este apartado, debe convertirse en el resumen de la experiencia práctica del proyecto, y por sí mismo justifica que la memoria se convierta en un documento útil, fuente de referencia para los autores, los tutores y futuros alumnos.
\section{Algoritmos para el cálculo de rutas}
Una de las decisiones más importantes del proyecto es la de elegir que algoritmo implementar, ya que es el pilar central de la ampliación que se quiere realizar.

El problema se trata de una variante del algoritmo base, que es el \textit{problema de la orientación} (\textit{OP}), detallado en la sección \textit{Conceptos teóricos}, concretamente la variante \textit{problema de la orientación con ventanas de tiempo} (\textit{OPTW}).

Las bases del algoritmo son sencillas: se quiere visitar una ciudad y se han de seleccionar que puntos de interés (ordenados) crean la ruta más óptima. No suele darse en caso en que se seleccionen todos los puntos disponibles, pero esto es por las restricciones de tiempo.

La diferencia principal con el algoritmo base es que el tiempo de viaje entre puntos ya no es el único tiempo a tener en cuenta. A cada punto de interés se le asigna una ventana de tiempo en la que será obligatorio llegar para que el punto sea visitable. Se puede pensar que la solución a este problema es tan sencilla como añadir una condición al proceso de selección, pero la verdadera complicación es cómo reducir lo máximo posible los tiempos de cómputo.

La decisión de que algoritmo de los disponibles implementar se basó en su extensibilidad, es decir, debido a que OPTW no es la única variación del algoritmo base, se eligió una solución que permitiese, con no demasiados cambios, resolver alguna de dichas variaciones en futuros proyectos. Dicho algoritmo el el basado en \textit{colonia de hormigas}.

Una de las desventajas que se encontraron con ese algoritmo fue la falta de tablas en las que se mostrasen comparaciones de los tiempos de cómputo, aunque ya se adelantaba que eran altos.
Es por eso que se decidió implementar un segundo algoritmo, a poder ser basado en una resolución lo más diferente posible, para así poder comparar como de eficiente era cada uno. Este segundo algoritmo está basado en \textit{búsquedas iterativas}.

Los dos algoritmos solución elegidos fueron:
\begin{itemize}
\tightlist
\item Iterated local search for OPTW
\item Ant Colony System
\end{itemize}
Las implementaciones de las dos soluciones se encuentran explicadas detalladamente en el \textit{Anexo IV: 2.Algoritmos para el cálculo de rutas}.

\section{Obtención y tratamiento de datos}

\subsection{Obtención de horarios}
Se evaluaron dos posibles fuentes de datos:
\begin{itemize}
\item Google Maps
\item OpenStreetMaps
\end{itemize}

Lo primero que se investigó fue el formato de los datos en cada una de las opciones, centrándonos sobre todo en los horarios de los establecimientos ya que eran los nuevos datos a implementar.

Por un lado \textit{Google Maps} es uno de los servicios más utilizados cuando se requieren datos de geolocalización. Cuenta con una API\cite{wiki:googlemapsapi} relativamente sencilla de usar, ofrece unos datos en un único formato y que (normalmente) han sido verificados, lo cual la hace una elección muy llamativa.

Cada punto cuenta con 3 categorías y cada una contiene diferentes campos. Dichas categorías y datos son:

\textbf{Basic}
\begin{itemize}
\tightlist
\item address\_component
\item adr\_address
\item formatted\_address
\item geometry
\item icon
\item name
\item permanently\_closed
\item photo
\item place\_id
\item plus\_code
\item type
\item url
\item utc\_offset
\item vicinity
\end{itemize}

\textbf{Contact}
\begin{itemize}
\tightlist
\item formatted\_phone\_number
\item international\_phone\_number
\item opening\_hours
\item website
\end{itemize}

\textbf{Atmosphere}
\begin{itemize}
\tightlist
\item price\_level
\item rating
\item review
\item user\_rattings\_total
\end{itemize}

Este servicio no solo ofrecía los datos necesarios, sino que aportaba otros extra que podían ser utilizados para hacer la aplicación más realista, permitiendo por ejemplo descartar ciertos puntos de interés en función del precio.

Al ir a implementar la API se encontró un obstáculo insalvable, las licencias de Google Maps.
Si bien es cierto que la restricción del número de elementos por consulta (100 elementos) ha existido desde el nacimiento de esta API, a partir del 16 de Julio de 2018 es necesario aportar los datos de facturación en la cuenta con la que se utiliza dicho servicio, ya que en esta fecha se estableció un coste por el uso de los datos.

Los algoritmos utilizados necesitan una gran cantidad de datos, ya que para disminuir el tiempo de cómputo se utiliza una \textit{matriz de distancias}, en la que se cargan los datos de localización de cada punto para facilitar la obtención del espacio entre estos.

Aunque era posible solventar la restricción de peticiones disminuyendo considerablemente los datos disponibles para los algoritmos, no había opción con el pago por el servicio, por lo que se tuvo que descartar la opción de utilizar Google Maps.

La segunda opción disponible era utilizar el proyecto colaborativo \textit{OpenStreetMaps}\cite{wiki:openstreetmaps}, que contaba con la ventaja de tener los datos abiertos al publico, ya que es la comunidad quien se encarga de aportar dichos datos. Para ello el único requisito es estar registrado.

Este último punto plantea otro inconveniente, dado que cualquiera puede añadir datos nuevos (o corregir los antiguos), era posible encontrarnos con ciertas incongruencias en los datos, sobre todo en los referentes a horarios de apertura y cierre de los establecimientos.

Otros inconvenientes de utilizar OpenStreetMaps eran:
\begin{itemize}
\tightlist
\item Falta de una API
\item Documentación oficial escasa del uso de sus herramientas
\item Necesidad de tener los datos previamente volcados en nuestro servidor para poder utilizarlos
\item Necesidad de una base de datos GIS
\end{itemize}

Dado que no había otra alternativa, la elección estaba tomada sin importar los inconvenientes. 

Cabe destacar el gran trabajo que hicieron los anteriores contribuyentes de este proyecto: Íñigo Vázquez Gómez, Roberto Villuela Uzquiza y Alejandro Cuevas Álvarez. Su trabajo facilitó mucho la comprensión y el uso de esta herramienta

\subsection{Volcado de datos, Osmosis vs Osm2pgsql}
Se investigó que herramientas existían para el volcado de los datos presentes en los ficheros descargados de \textit{OpenStreetMaps} a nuestra base de datos geoespacial. Se valoraron dos alternativas:
\begin{itemize}
\item Osmosis
\item Osm2pgsql
\end{itemize}

Una de las ventajas de  Osmosis es que los datos se vuelcan de una forma muy particular, creando relaciones de las que podíamos beneficiarnos como ya habían hecho en años anteriores el resto de contribuyentes de esta aplicación.

Una vez volcados los datos se reparó en que no solo no se había importado la etiqueta opening hours, sino que no se había importado ninguna etiqueta que no fuese de las principales, entendiendo por principales:
\begin{itemize}
\item ID
\item Nombre
\item Localización
\item Tipo de nodo/vía
\item Relaciones entre nodos
\end{itemize}

En vista del resultado se pasó a utilizar la herramienta \textit{Osm2pgsql} ya que consta de un archivo de estilo, llamado \textit{default style}, en el que se pueden indicar que etiquetas de todas las existentes en OpenStreetMaps queremos importar.
El volcado se realizó con éxito, pero el formato de las tablas creadas en \textit{PostgreSQL} no eran adecuadas para el uso que se les quería dar.

Finalmente, se decidió que, debido a que cada herramienta aportaba utilidades imprescindibles, se utilizarían ambas herramientas, beneficiándonos así de sus puntos fuertes y solventando las debilidades.
La solución final se encuentra explicada detalladamente en el \textit{Anexo RELLENAR ANEXO}.

Cabe destacar que en caso de querer volcar una gran cantidad de datos, como por ejemplo, de toda España, existe la posibilidad de modificar la cantidad de memoria de nuestro PC que permitimos utilizar a Osm2pgsql. Si dicha cantidad no es suficiente, a mitad del volcado se quedará sin espacio y se nos notificará con un error que no se han podido volcar los datos. Por el contrario, si ofrecemos demasiada memoria nuestro ordenador dejará de responder y tendremos que forzar un apagado del sistema manteniendo pulsado el botón de encendido.

Se recomienda que, en caso de necesitar utilizar ficheros tan pesados, es mejor obtener cortes de dicho fichero, por ejemplo, por comunidades autónomas o por provincias y realizar el volcado de uno en uno, para así evitar errores como los mencionados.

\subsection{Tratamiento de horarios}
Tras el volcado de datos se pasó a comprobar la calidad de la información.

La etiqueta \textit{opening\_hours} cuenta con una sintaxis bastante compleja, que bien utilizada es capaz de aportar una gran información. Un ejemplo de dicha sintaxis es el siguiente:
\imagen{opHoursExample}{Ejemplo de sintaxis de la etiqueta opening hours \cite{wiki:opHours}}

Como se ha comentado con anterioridad, OpenStreetMaps permite a sus usuarios subir sus propios datos, lo que causa que algunos de ellos contengan una sintaxis errónea.

Después de comentar con el tutor del proyecto las posibles soluciones, se llegó a dos ideas:
\begin{itemize}
\item Gestión de los datos desde la base de datos
\item Gestión de los datos desde los algoritmos
\end{itemize}

Las comprobaciones de los datos mostraban que una gran cantidad de puntos de interés no contaban con un horario, pero no todos los usuarios habían optado por no introducir ningún valor, ciertos establecimientos tenían en su horario ``", lo cual no era detectado ni como nulo ni como espacio. Esto dio varios dolores de cabeza a la hora de probar los algoritmos hasta que se descubrió que esta era la causa.

Se decidió por tanto practicar las dos soluciones a la par, nada más introducir los datos se usaría un script SQL que buscaría aquellos que contuviesen tanto nulo como ``" y sustituiría esos valores por unos por defecto.

Por otro lado, se estudiaron todos los posibles tipos de sintaxis y se creó un método dentro del algoritmo que extrajese los horarios de apertura y cierre a partir de un String obtenido de la base de datos.

Estas soluciones se encuentran explicadas mas detalladamente en el \textit{Anexo RELLENAR ANEXO}

\section{Conocimientos adquiridos en la carrera}
\subsection{Conocimientos sobre inteligencia artificial}
Ha sido de gran ayuda para este proyecto haber cursado durante la carrera la asignatura \textit{Computación Neuronal y Evolutiva}, ya que tuvimos nuestro primer acercamiento a este tipo de algoritmos mediante un desafio Hash Code, así como las asignaturas \textit{Sistemas Inteligentes}, \textit{Algoritmia}, la asignatura \textit{Advanced Logical and Functional Programming} cursada en la \textit{West University of Timisoara} etc.
Gracias a esto la comprensión de artículos referentes a la resolución e implementación del problema de la orientación ha sido mucho mas llevadera.

\subsection{Conocimientos sobre bases de datos}
De vital importancia ya que ha sido necesario crear una base de datos GIS y modificarla mediante SQL para que fuese más eficiente.
Gracias a haber cursado las asignaturas de \textit{Bases de Datos y Aplicaciones de Bases de Datos} se ha partido de una posición más avanzada y ha sido posible tanto la creación de la base de datos como el manejo del álgebra relacional.

\subsection{Trabajo de investigación}
No se ha impartido en ninguna asignatura en concreto, pero ha estado presente en cada uno de los trabajos que se han desarrollado en el grado.

La rapidez con la que se buscaban y comprendían nuevos conocimientos, se verificaban las fuentes de los datos y se solventaban los diversos problemas y errores que han aparecido en el desarrollo de este proyecto ha demostrado que esta capacidad es una de las más importantes que se han adquirido durante la carrera. 

\capitulo{6}{Trabajos relacionados}

Este apartado sería parecido a un estado del arte de una tesis o tesina. En un trabajo final grado no parece obligada su presencia, aunque se puede dejar a juicio del tutor el incluir un pequeño resumen comentado de los trabajos y proyectos ya realizados en el campo del proyecto en curso. 

\capitulo{7}{Conclusiones y Líneas de trabajo futuras}
En este apartado se tratarán las diferentes conclusiones que han surgido durante la realización de este proyecto así como posibles lineas que seguir en un futuro para mejorar la funcionalidad de la aplicación.

\section{Conclusiones}
Durante la realización de este proyecto se ha aprendido a utilizar repositorios como OpenStreetMaps para obtener datos geográficos, a crear y utilizar bases de datos geoespaciales, a trabajar con proyectos Java bajo la herramienta de gestión Maven, a implementar y comprender diferentes algoritmos para resolver el problema de la orientación, a trabajar con el servidor de aplicaciones Glassfish y a tratar con los diversos problemas que genera el trabajar en la ampliación de un proyecto anterior.

Como valoración personal, con este proyecto se han conseguido reforzar diferentes conocimientos adquiridos a lo largo de toda la carrera, muchos de los cuales no se habían puesto en práctica.

\section{Líneas de trabajo futuras}
En este apartado se comentarán las posibles mejoras que se podrían aplicar a este proyecto:
\begin{itemize}
\item Ampliar la generación de rutas, haciendo al algoritmo dependiente del tiempo, es decir, que la velocidad de ruta varíe en función de la hora del día (Time Dependent Orienteering Problem). Los algoritmos implementados en este proyecto están preparados para dicha ampliación.
\item Crear una página web que iguale la funcionalidad de la aplicación.
\item Conectar la aplicación con las redes sociales, para que permita compartir y publicar las rutas generadas.
\item Conectar la aplicación con el GPS del dispositivo, para que pueda guiarnos hacia el siguiente punto de interés.
\item Actualizar el cliente a una versión de Android más actual.
\item Crear una versión del cliente para otros sistema móviles como pueden ser IOS o Windows Phone.
\item Incluir realidad aumentada en la aplicación, para poder ver los diferentes puntos de interés sin necesidad de hacer la ruta personalmente.
\item Conectar la aplicación a diferentes funciones del dispositivo para, en función de la distancia, el tiempo recorrido y los pasos dados, proveer al usuario información sobre su salud física.
\item Un sistema de amigos con el que se puedan ver las últimas rutas de estos.
\end{itemize}


\bibliographystyle{plain}
\bibliography{bibliografia}

\end{document}
